\documentclass[11pt]{article}
\usepackage[margin=1in]{geometry}
\usepackage{booktabs,hyperref,graphicx}
\title{GuardSynth: Verified Constraint Rendering for Visual Action Learning}
\author{Author information pending}
\date{Working manuscript, September 2026}
\begin{document}
\maketitle
\begin{abstract}
We study whether a source-bearing executable constraint representation can preserve
previously observed benefits of supplied natural-language constraints in visual
language model action selection. Our pipeline separates supplied assumptions,
typed constraint lowering, bounded solver checks, controlled-language rendering,
and visual model learning. A controlled temporal task compares an unconstrained
chain of cognition (P0), a supplied natural-language constraint (P1), and an
EBLC-Core-verified rendering of the same constraint (P2). Constraints are visible
at training and inference in both constrained conditions. Nine local LoRA fits
completed 648 optimizer updates and 2,592 action predictions. On the standard
benchmark, P0 violation/safe completion was 0.319/0.681; both P1 and P2 achieved
0.000/1.000 across all three seeds. A new generator-seed sample reproduced this
pattern. The prospectively fixed point criteria support preservation within this
small reused template family, not formal noninferiority or road generalization.
\end{abstract}
\section{Question and scope}
Verification of a formal representation does not establish the truth of its
source assumptions or the behavior of a learned policy. We therefore distinguish
semantic correspondence, behavioral effect preservation, and real-scene source
applicability. The present experiment tests the first two under supplied synthetic
assumptions. Existing reviewed road scenes provide a separate application case
study, not an independent road-safety test.
\section{Attributed prior evidence}
Project01's temporal guard experiment trained Qwen3-VL-2B LoRA adapters for three
seeds. Its reported standard-task mean violation rates were 0.319 for original
CoC and 0.000 for inline natural-language constraints; safe goal completion was
0.681 and 1.000. These are prior Project01 results, not current P2 measurements.
On unseen time values, safe completion was 0.625 and 0.649, indicating limited
numerical generalization. We reuse that project's generator, visual processor,
and optimizer harness without modifying its source or results.
The historical task contains deliberately supplied guards rather than automatically
extracted road constraints. Its earlier inline-NL condition remains P1's reference;
it is never relabeled as a previously demonstrated P2 result.
\section{Method}
The restricted supplied temporal profile identifies ego, a pedestrian conflict
zone, a no-entry obligation, and a continuous-clear release duration. The world
assumes exactly one occupied-to-clear transition without reoccupation. Lowering
uses integer millisecond tick counts and an implication from entry to entry time
being no earlier than clearance plus the duration. The shared EBLC Core parser
checks typing and the SMT compiler checks feasibility and boundary cases.
The renderer reads only validated contract fields; candidate labels and action
answers are not inputs to constraint generation. This is a controlled profile,
not a claim that the general EBLC lifecycle or real-road extraction is solved.
The modeled decision is entry admissibility: candidate waiting periods are fixed
by the controlled task. The Core does not model velocity during a wait; the CNL
comparison does not certify kinematic equivalence outside this candidate family.
\section{Frozen experiment}
Protocol \texttt{paper1-method-preservation-v0.1} fixes seeds 42, 17, and 123;
192 training scenes, two counterfactual contracts per scene; three epochs,
batch size two, accumulation eight, and 72 optimizer updates per arm and seed.
All arms use the same Qwen checkpoint, visual inputs, targets, and optimizer.
P1 and P2 append their constraint to the original CoC during both learning and
inference. Only assistant answer tokens receive loss. P0 deliberately lacks the
contract needed to distinguish counterfactual pairs.
\begin{figure}[t]
\centering
\includegraphics[width=\linewidth]{artifacts/projects/guardsynth-coc/public/paper1-method-preservation-001/temporal-paper-2026-09-29-001/storyboards.pdf}
\caption{All six distinct rendered storyboards in the reused Project01 task.
Each contains four temporal panels. Repeated generator seeds change sample
frequencies and candidate label rotations, not this visual environment family.}
\end{figure}

Each checkpoint is evaluated on the historical benchmark distribution, a new
fixed generator-seed sample, and unseen temporal values. The independent program
scorer uses environment times and candidate progress, not generated CNL or model
self-assessment. We report violations, safe goal completion, deadlock as an
unnecessary-stop proxy, accuracy, coverage, paired contract accuracy, and raw
per-example outputs. Template-cluster paired bootstrap intervals describe
uncertainty conditional on this small task family.

Before new outcomes, we fixed per-seed level criteria: P1 and P2 violation at
most 5\% and safe completion at least 95\%; P2 improvement over P0 at least
20 percentage points on both violation and safe completion; and no more than
5 percentage points degradation from P1 on either endpoint. The five-percentage-point tolerance is a prospectively chosen operational
criterion, not a margin supplied by the historical study. These operational
criteria do not establish formal noninferiority. Equally weak P1 and P2 fail.
\section{Results}
All nine fits completed 72 updates each and saved distinct adapters. Each adapter
contains 224 tensors, including 112 nonzero LoRA B matrices. All 2,592 predictions
were valid. Table~\ref{tab:results} reports the new Project04 measurements;
historical values were not copied into these cells. P0/P1 deliberately reuse the
prior generator, seeds, and task inputs as replication controls, not independent
new training cohorts. P2 uses a different rendered constraint string.
\begin{table}[ht]
\centering
\begin{tabular}{llrrr}
\toprule
Evaluation & Arm & Violation & Safe completion & Accuracy\\
\midrule
Benchmark & P0 & .3194 & .6806 & .5000\\
          & P1 & .0000 & 1.0000 & 1.0000\\
          & P2 & .0000 & 1.0000 & 1.0000\\
New seed  & P0 & .3160 & .6840 & .5000\\
          & P1 & .0000 & 1.0000 & 1.0000\\
          & P2 & .0000 & 1.0000 & 1.0000\\
Unseen time & P0 & .3750 & .6250 & .5000\\
            & P1 & .3611 & .6389 & .6250\\
            & P2 & .0000 & 1.0000 & .8993\\
\bottomrule
\end{tabular}
\caption{Three-seed means from newly executed matched fits. Deadlock was zero
and coverage was one for all arms/splits. New seeds reuse the same templates.}
\label{tab:results}
\end{table}

Every seed met the prespecified P1 replication, P2 level, P2--P0 improvement, and
P2--P1 preservation point criteria on both standard evaluations. The new-seed
P2--P0 violation difference was $-0.3160$, with conditional template-cluster
95\% bootstrap interval $[-0.4348,-0.1832]$; safe completion increased by the
same magnitude. P2--P1 differences were zero, with degenerate empirical intervals
$[0,0]$. These intervals describe the observed templates and do not prove a
population noninferiority margin. In particular, the interval does not establish
a minimum 20-percentage-point improvement over broader environments.

The unseen-time results are descriptive secondary evidence: P1 safe completion
remained limited at .6389, whereas P2 reached 1.0000 with accuracy .8993. The
gap between completion and accuracy reflects safe but suboptimal later entries.
Contract-pair accuracy was .0000/.2778/.7986 for P0/P1/P2 on unseen time, versus
.0000/1.0000/1.0000 on the standard evaluations. This does not isolate a solver
benefit or establish general numerical reasoning.

The bridge passes 100 SMT queries covering
four release durations and boundary/holding cases. The supplied contract bridge also rejects unsupported subject, unit, obligation,
and duration changes. A deleted-entry-gate mutation is detected by the boundary
queries. The paired data export has
2,016 records across seed/split combinations; these are not 2,016 independent
visual scenes.
\section{Limitations}
The historical renderer generates only three distinct standard storyboards and
three unseen-time storyboards. New generator seeds reuse those templates and
candidate rotations; they are not fresh visual environments. The test exposes
contract information at inference and cannot support a claim about inference
without constraints. Supplied synthetic rules are not automatically extracted
road facts. Solver satisfiability does not prove learned safety. Three model
seeds do not replace independent environments. Negative or inconclusive
preservation results remain reportable.
Actual image tensors enter every model call, but this alone does not demonstrate
visual grounding: the regular candidate time schedule can also reveal clearance
timing. No image-ablation or independent natural-image transfer claim is made.
The P1--P2 comparison evaluates preservation through the verified rendering
pipeline, not the causal benefit of running a solver. Wording and representation
change together; a verification-omission control would be required for that
additional attribution claim. A successful result does not establish superiority
of EBLC over every simpler representation.
The comparison also changes constraint exposure jointly during training and
inference. Actual optimizer updates are verified, but their causal contribution
cannot be separated from inference-time conditioning without an additional
matched pretrained control in this new protocol.
\section{Real-scene case study}
Existing single-reviewer development ACTION references, source confirmations,
unknown-motion states, scoped exclusions, and disagreements are preserved.
They are not consensus gold. The two accepted STOP cases and the separately
reviewed scene18 launch check do not establish road-level effectiveness.
The frozen pretrained Qwen development baseline produced 19 valid predictions
from 133 causal frames on 17 previously exposed clips, with no optimizer updates.
Against the preserved scoped ACTION references, inappropriate selections were
10/16, normal-progress selections 2/8, and unnecessary stops 7/12. Their distinct
denominators reflect unknown references and three clarified field masks. These
appropriateness assessments are not certified collision or safety violations.
The separate scene18 launch check saved one-update adapters for two conditions;
it is execution evidence rather than efficacy evidence. No road review was
repeated or generated to support the controlled temporal experiment.
\section{Reproducibility}
The owner of new code and results is Project04, \texttt{guardsynth-coc}.
The paired protocol export pins the original Project01 world and trainer hashes,
six original result records, the renderer, and the fixed protocol. Original runs
remain immutable. New adapters, losses, raw per-example predictions, and output
hashes are stored under the owner-scoped \texttt{paper1-method-preservation-001}
artifact namespace. The run uses only local resources and the pinned Qwen
revision:
\begin{quote}\small\texttt{89644892e4d85e24eaac8bacfd4f463576704203}\end{quote}
No expanded LLM labels were needed for this supplied-rule task; consequently,
no human sample-audit response or machine label is represented as human gold.
\end{document}
